\section{Corroborate}
\label{sec:corroborate}

\judge{Corroborate} builds its reference from the one resource every dataset has,
the responses themselves. It needs no LLM, no document collection and no network
access.

\subsection{Pooling sentences into claims}

For each topic we take every sentence of every response and normalize it. We drop
fixed boilerplate and any sentence shorter than five words. Each sentence is
embedded with bge-base-en-v1.5~\cite{xiao2023cpack}, a model that maps sentences
with similar meaning to nearby vectors. We then cluster them greedily, in a fixed
order by run id and sentence position. A sentence joins the cluster whose first
sentence, the leader, is most similar to it, as long as the cosine similarity is at
least $\theta=0.80$. Otherwise it starts a new cluster. The fixed order makes the
result independent of file order. Each cluster records the set of teams $T(c)$
whose responses contributed to it.

Clustering also does the matching. A response covers a cluster when one of its own
sentences ended up in it, so we need no separate retrieval step.

Both $\theta$ and the partial-credit cutoff of 0.62 used below come from similarities on a small synthetic dataset. It ships with the task's starter kit, the template code participants build their judges from. There a true paraphrase scored about 0.83
while an unrelated sentence on the same topic got about 0.61. Neither was swept on
the pilots.

\subsection{The claim bank}

A cluster enters the topic's bank only if at least two teams contributed to it. It
must also be specific to the topic. Specificity is the cosine similarity between
the cluster's leader and the topic text, with a floor of 0.50. The bank keeps at
most the 40 clusters with the most teams. Votes are counted per team, never per
run. Otherwise a team that submits five variants of one pipeline would
vote five times for its own phrasing. We set the cap of 40 and the five-word
minimum by hand and did not tune them.

\subsection{Leave-one-team-out weighting}

A response's own sentences are part of the pool behind its bank. Two
rules keep a team from voting for its own claims. First, a cluster needs at least
two teams to enter the bank. So a claim made by only one team is never scored.
Second, when we score a response from team $t$, each cluster gets the weight
\begin{equation}
  w_t(c) = |T(c) \setminus \{t\}|,
\end{equation}
the number of other teams that state it. Without this weight, a team's own
sentences would raise the weight of the claims it is scored on. Every bank
cluster has at least two teams, so the weight is never below one. Leaving the
scored team out of every weight is what we mean by leave-one-team-out.

The scheme is a team-level version of the leave-one-run-out test
of~\citet{zobel1998reliable}. That test removes from the pool whatever one run
alone contributed. Zobel used it to measure pool bias. We use the same idea to
correct for it.

\subsection{Measures}

Coverage is soft and allows partial credit. For a bank cluster $c$, response $r$
gets $\mathrm{cov}_r(c)=1$ if one of its sentences was assigned to $c$. The same
holds if its most similar sentence reaches $\theta$. Below 0.62 the credit is 0. In
between, when the best similarity falls in $[0.62, \theta)$, the credit is the
entailment probability from a cross-encoder. A cross-encoder reads both texts
together and scores the pair directly, instead of comparing separate embeddings.
Ours is nli-deberta-v3-base, a DeBERTaV3 model~\cite{he2023debertav3} fine-tuned
for natural language inference. That is the task of deciding whether one text
entails, contradicts or is neutral to another. Coverage is then
\begin{equation}
  \texttt{NCOV}(r) = \frac{\sum_{c} w_t(c)\,\mathrm{cov}_r(c)}{\sum_{c} w_t(c)}.
\end{equation}

\texttt{NPREC} is the share of the response's own sentences that sit in a
topic-specific cluster also stated by at least one other team. Because it is
normalized by length, padding and claims nobody else makes pull it down.
\texttt{NF1} is the harmonic mean of \texttt{NCOV} and \texttt{NPREC}.
\texttt{NF1\_RANK} is the within-topic percentile of \texttt{NF1}, as in
Equation~\ref{eq:percentile}. 

\texttt{CITE} measures citation support. For each cited sentence we check the first
three documents it cites. We split each document into chunks of at most 1,200
characters and use only its first six. For each document we take the chunk closest
to the sentence by embedding similarity. The cross-encoder's probability that this
chunk entails the sentence is the document's support. A sentence's support is the
highest over its documents. \texttt{CITE} averages it over the cited sentences that
have at least one document present. The path is chosen per dataset. When fewer
than half the responses carry document text, \texttt{CITE} uses citation overlap
instead. This is the mean share of other teams citing the same documents, times the
share of sentences that cite anything. Both test sets carry the text, so on them
\texttt{CITE} uses entailment.

\subsection{Why a pooled bank is not circular}

Two arguments support the design. First, the human reference is itself built from
a pool. In TREC RAG, the documents that systems retrieved are pooled and nugget
banks are drafted from them. Assessors then vet the
banks~\cite{pradeep2024initial}. That pool holds documents rather than answers, so
the analogy is loose. Still, in both cases the reference comes from what the
systems produced. Second, a fact that no system states cannot enter the bank. It
would give every system the same zero anyway, so it cannot change their order.

The real risk is that consensus pulls every system toward the median. That would
hurt most at the top of the ranking, where \tauap\ weighs errors most.
Section~\ref{sec:analysis} suggests this is what happens on short answers. We have
not tested it directly.

\subsection{Cost}

Our first version issued one matrix-vector product per sentence. Nearly all of its
time went to starting and joining thread pools. We now cluster in blocks of 512
sentences, which cut pooling from 13.0 to 1.04 seconds per topic. A test checks
that the output is identical to the sequential version. On one GPU, a full run
takes 73 minutes on rag26 and 22 minutes on ragtime26. On CPU alone, a pilot took
about eight hours.
